% Job posting: https://huaweicanada.recruitee.com/o/senior-researcher-edge-ai-optimizationhardware-aware-ml
\documentclass[11pt]{article}
\usepackage[left=0.7in,top=0.7in,right=0.7in,bottom=0.7in]{geometry}
\usepackage{enumitem}
\usepackage[hidelinks]{hyperref}
\usepackage{titlesec}
\usepackage{array}
\usepackage{setspace}
\usepackage{needspace}
\usepackage[T1]{fontenc}
\usepackage{newtxtext,newtxmath}
\ifdefined\pdfgentounicode
  \input{glyphtounicode}
  \pdfgentounicode=1
\fi
\setstretch{1.04}
\pagenumbering{gobble}
\titleformat{\section}{\Large\scshape}{}{4em}{}[\vspace{-0.7em}\rule{\linewidth}{0.4pt}\vspace{-0.4em}]
\titlespacing*{\section}{0pt}{1em}{0.4em}
\newenvironment{cvrole}[5]{%
  \par\noindent\textbf{\large #1} | \textit{\small #2, #3}\hfill \textit{\small #4 - #5}\par
  \begin{itemize}[leftmargin=2em, itemsep=-0.3em, topsep=-0.5em]%
}{%
  \end{itemize}\par\noindent\mbox{}\par\vspace{0.1em}%
}
\newcommand{\cvName}{Nishal Stanislaus Silva}
\newcommand{\cvEmail}{nishal.silva@hotmail.com}
\newcommand{\cvPhone}{+1 613 200 2030}
\newcommand{\cvCity}{Toronto, ON}
\newcommand{\cvSite}{https://nishal.xyz}
\newcommand{\cvLinkedIn}{https://www.linkedin.com/in/nishal-silva}
\newcommand{\cvGitHub}{https://github.com/ns2max}
\newcommand{\cvStatus}{Canadian Permanent Resident, Eligible to work in Canada without sponsorship}
\newcommand{\cvTagline}{ML Research Engineer | Real-Time \& Edge Inference | Audio, Sensor \& Time-Series}
\begin{document}
\pagestyle{empty}
\begin{center}
    {\LARGE \scshape \textbf{\cvName}}{\large \textit{, PhD}}\\[1pt]
    {\small \cvTagline}\\[1pt]
    {\footnotesize\textit{\cvStatus}}\\[1pt]
    {\small
    \href{mailto:\cvEmail}{\cvEmail} \,|\, \cvPhone \,|\, \cvCity
    \,|\, \href{\cvSite}{nishal.xyz} \,|\, \href{\cvLinkedIn}{linkedin.com/in/nishal-silva}
    \,|\, \href{\cvGitHub}{github.com/ns2max}}
\end{center}

\section*{Professional Summary}
PhD researcher whose doctoral work is, at its core, hardware-aware ML: designing and validating models under a fixed latency, compute, and power envelope on embedded ARM hardware. Headline result: a real-time audio pattern detector running at 14 ms inference on a Raspberry Pi 4, F1 0.76, 31.4\% CPU utilization -- 74x faster than the 1038 ms DTW baseline it replaced (JAES 2026). Built MIRaaS, a UDP edge-to-server offload architecture that characterizes exactly where accuracy, latency, and compute trade off when deciding what runs on-edge versus off-edge (accepted, IEEE IS2 2026). Track record of frozen-protocol benchmark design, ablation studies, and quantified accuracy/latency/CPU trade-off reporting across five years of applied ML on embedded and industrial hardware.

\section*{Core Competencies}
\textbf{Hardware-Aware ML \& Edge Optimization:} ARM CPU and power profiling; real-time inference under sub-30 ms budgets; embedded Linux deployment (Raspberry Pi 4); model compression, quantization, and ONNX/TFLite familiarity; frozen-protocol benchmark and ablation design; latency/accuracy/CPU trade-off quantification; edge-vs-server offload architecture. \textbf{Also:} real-time systems engineering, sensor/time-series pattern detection, computer vision at manufacturing scale, ML pipeline ownership (acquisition through deployment and monitoring).

\section*{Technical Stack}
C++17, Python, TensorFlow, Keras, PyTorch, scikit-learn, ONNX, TFLite (familiarity), embedded Linux, Raspberry Pi/ARM, Docker, CI/CD, AWS, SQL, ETL pipelines, Git, pytest.

\section*{Education}
\textbf{PhD, Information Engineering and Computer Science} -- University of Trento, Italy \hfill 2021 - 2025 \\
\textit{Thesis: "Embedded Real-time Musical Pattern Detection for Smart Musical Instruments" -- hardware-aware model design under fixed latency/compute/power budgets on embedded ARM hardware.} \\
\textbf{MSc, Telecommunication and Electronic Engineering} -- Sheffield Hallam University, UK \hfill 2016 - 2018 \\
\textbf{BEng (Hons), Electronic Engineering} -- Sheffield Hallam University, UK \hfill 2013 - 2014

\section*{Research and Work Experience}

\begin{cvrole}{Independent Researcher}{Self-directed}{Toronto, Canada}{Jan 2026}{Present}
\item Authored two accepted peer-reviewed papers (Smart Drums, JAES; MIRaaS, IEEE IS2 2026) continuing the hardware-aware ML and edge-vs-server characterization line of work.
\item Maintain Nebula, a C++17 audio-feature-extraction library with per-feature ARM/Raspberry Pi latency benchmarks (github.com/ns2max/nebula, 13 stars).
\item Use AI-assisted development workflows, including Claude Code, to accelerate library and tooling development.
\end{cvrole}

\begin{cvrole}{Postdoctoral Researcher}{University of Trento}{Trento, Italy}{Jan 2025}{Dec 2025}
\item Owned the full ML pipeline for streaming audio and IMU/gesture data on the EU Horizon EIC Pathfinder project MUSMET: acquisition, DSP/feature engineering, training and evaluation, edge deployment on Raspberry Pi 4, and monitoring.
\item Designed and shipped a real-time audio pattern detector at 14 ms inference, F1 0.76, 31.4\% CPU on Raspberry Pi 4 -- 74x faster than the 1038 ms DTW baseline it replaced (JAES 2026).
\item Designed frozen-protocol benchmark suites and ablations (RNN vs.\ DTW; audio vs.\ MIDI; pattern-set sizes 1/3/10) quantifying accuracy, latency, and CPU trade-offs, with a written record of negative results.
\item Built MIRaaS, a UDP edge-to-server offload architecture characterizing latency/compute trade-offs for on-edge vs.\ off-edge inference decisions (accepted, IEEE IS2 2026).
\end{cvrole}

\begin{cvrole}{Visiting Researcher}{McGill University (IDMIL / CIRMMT)}{Montreal, Canada}{May 2024}{Aug 2024}
\item Built a self-powered smart electric guitar (IMU + audio) with gesture fusion via a hierarchical FSM, reaching F1 0.78 on a 500-event corpus (IEEE IS2 2025).
\item Profiled compute and power draw to validate real-time feasibility on self-powered embedded hardware -- power-budget hardware-aware engineering.
\end{cvrole}

\begin{cvrole}{Visiting Researcher}{University of Visual and Performing Arts}{Colombo, Sri Lanka}{Jan 2024}{Mar 2024}
\item Conducted a human-centred evaluation of musicians' perception of smart instruments with embedded pattern detection (IEEE I3DA 2025).
\end{cvrole}

\begin{cvrole}{Technical Consultant}{Forestpin}{Colombo, Sri Lanka}{Aug 2020}{Feb 2021}
\item Built ML and statistical pipelines for financial forensics, compliance monitoring, and risk detection on large structured enterprise data.
\end{cvrole}

\begin{cvrole}{Research Engineer}{MAS Holdings}{Colombo, Sri Lanka}{Jan 2015}{Jul 2020}
\item Delivered end-to-end computer-vision and IoT systems for manufacturing at scale over 5.5 years, raising operational efficiency by up to 300\%.
\item Ran real-time IoT ETL pipelines across sites, deploying C++/Python inference on embedded hardware, and introduced CI/CD and code review practices.
\item Automated multilingual care-label quality control (OpenCV, conveyor + PLC reject integration), cutting inspection time by 99.5\%.
\end{cvrole}

\section*{Publications}
Silva, N. \& Turchet, L. (2026). Real-Time Audio Pattern Detection for Smart Musical Instruments. \textit{J. Audio Eng. Soc.} 74(3). Silva, N. \& Turchet, L. (2026, in press). Smart Drums. \textit{J. Audio Eng. Soc.} Silva, N. \& Turchet, L. (2026, accepted). MIRaaS: Latency Characterization of an Edge-to-Server Architecture. \textit{IEEE IS2}. Silva, N. \& Turchet, L. (2022). A Training-Free SSIM-Based Method for Onset Detection, achieving 95\% detection with no training data. \textit{DAFx}.

\section*{Highlights}
Finalist, MIDI Innovation Awards 2023 (Prototypes \& Non-Commercial Software), for a real-time smart-instrument prototype; recognized by the Government Medical Officers' Association of Sri Lanka (2020) for a COVID-19 remote vitals-monitoring deployment; peer reviewer for IEEE Access and the Journal of the National Science Foundation of Sri Lanka; programme committee member for IEEE IS2 and IEEE I3DA; supervised two BSc students.

\end{document}
